\section{Related Work and Gap}
\label{sec:related}

CVSS is the reference scoring vocabulary for many vulnerability-management workflows. CVSS v3.1 organizes metrics into Base, Temporal, and Environmental groups, while CVSS v4.0 reorganizes the model into Base, Threat, Environmental, and Supplemental groups. This distinction is important for the present work because Base metrics describe intrinsic vulnerability properties, whereas Environmental metrics describe characteristics that are unique to the consumer's deployment environment \cite{firstcvss31,firstcvss40}. The official CVSS v4.0 specification explicitly recommends enriching Base metrics with Threat and Environmental values when the consumer needs a score that better reflects organizational risk context \cite{firstcvss40}.

Public vulnerability data sources do not usually solve this environmental-assessment problem for the consumer. The NVD provides CVSS enrichment for published CVE records, but its assessment is primarily based on Base metrics and it does not currently provide Environmental, Threat, or Supplemental metric assessments for a specific consuming organization \cite{nvdmetrics}. This creates the operational gap addressed by this paper: the information needed for environmental scoring is normally local to the organization, distributed across asset inventory, topology, firewall rules, business impact, and compliance scope.

Alternative prioritization frameworks address related but different problems. SSVC structures vulnerability response as a stakeholder-specific decision process, reducing reliance on one-size-fits-all scoring and emphasizing decisions that match the stakeholder's context \cite{ssvc2021,cisassvc}. EPSS estimates the probability that a vulnerability will be exploited in the wild, providing a threat-likelihood signal that complements severity scoring \cite{epss2021,firstepss}. These approaches are useful for prioritization, but they do not by themselves derive CVSS Environmental metrics such as security requirements or modified attack vector from local network evidence.

Prior environmental-scoring work is closest to the present study. Eitel proposed Environmental Aware Vulnerability Scoring, arguing that Environmental metrics are frequently ignored and that systematic automation requires information about systems and networks represented in a model \cite{eitel2020environmental}. The present work follows the same evidence-integration premise but uses the AI Bridge watcher as an execution and audit layer: it can read structured local artifacts, coordinate specialized review roles, regenerate outputs, and preserve an inspectable trace.

Recent work on language models and vulnerability scoring has focused mainly on automated CVSS prediction from vulnerability descriptions. AutoCVSS evaluates LLMs for automated software vulnerability scoring \cite{sanvito2025autocvss}, and EvalSVA studies multi-agent evaluators for software vulnerability assessment \cite{wen2025evalsva}. These systems are relevant because they show that LLMs and multi-agent designs can support vulnerability assessment tasks. However, their primary input is vulnerability text or code-related evidence, not organization-specific environmental artifacts. The remaining gap is therefore a reproducible workflow that connects local environmental evidence to CVSS Environmental metrics and records why each metric value was selected.
